% !TeX root = ../main.tex
\section{Methodology}
\label{sec:methodology}

\subsection{Hybrid Analog/Digital Deployment}
\label{subsec:hybrid-deployment}

We adopt a hybrid analog/digital inference stack: dense linear operators execute on differential-pair crossbars; control-heavy operators remain digital. Figure~\ref{fig:system-architecture} summarizes this division at the system level, and Table~\ref{tab:operator-mapping} lists the operator-level rationale following established practice in heterogeneous CIM accelerators \citep{wu2023bwq,wang2024epim,ge2024allspark,kim2025hemlet,lin2024hardsea}.

Under this mapping, 87.7\% of parameters are assigned to analog execution, whereas 57.9\% of the estimated energy remains in the digital domain, largely because of attention operations. This split reflects a basic architectural constraint: analog crossbars are well suited to weight-stationary matrix multiplication, but much less suited to dynamic, input-dependent computations \citep{kim2025hemlet}. Recent heterogeneous CIM accelerators for Vision Transformers follow the same logic, keeping attention digital even as linear projections move to analog arrays \citep{lin2024hardsea,wang2024epim,ge2024allspark}. The present organic device literature likewise demonstrates regular static array programming more naturally than fully dynamic interaction engines \citep{gebregiorgis2023organiccim,zhang2026opect}.

\begin{figure}[t]
    \centering
    \resizebox{0.98\textwidth}{!}{
    \begin{tikzpicture}[
        >=Latex,
        font=\small\rmfamily\upshape,
        node distance=6mm and 8mm,
        stage/.style={draw=black!60, rounded corners=4pt, thick, fill=black!2, minimum height=30mm, align=center},
        analog/.style={draw=teal!70!black, rounded corners=4pt, thick, fill=teal!8, minimum height=13mm, align=center},
        digital/.style={draw=black!55, rounded corners=4pt, thick, fill=black!5, minimum height=11mm, align=center},
        arrow/.style={->, thick, draw=black!65}
    ]
    \node[stage, minimum width=27mm] (input) {\textbf{Input}\\edge-vision image};
    \node[stage, minimum width=34mm, right=of input, draw=teal!70!black, fill=teal!10] (front) {\textbf{Physical frontend}\\photoresponse\\preprocessing};

    \node[draw=black!45, dashed, rounded corners=6pt, inner sep=6pt, fit={(0,0) (0,0)}] (dummy) {};
    \node[analog, minimum width=54mm, right=13mm of front, yshift=8mm] (analogpath) {\textbf{Analog CIM path}\\patch embedding, QKV / output, MLP};
    \node[digital, minimum width=54mm, below=7mm of analogpath] (digitalpath) {\textbf{Digital path}\\depthwise conv, attention, softmax, LN, head};
    \node[draw=black!45, dashed, rounded corners=6pt, inner sep=8pt, fit={(analogpath)(digitalpath)}] (backbone) {};
    \node[font=\bfseries, fill=white, inner sep=2pt, above=1mm of backbone.north] {Hybrid backbone};

    \node[stage, minimum width=38mm, right=11mm of analogpath, yshift=-8mm] (periph) {\textbf{Peripheral / calibration}\\ADC / DAC\\scale recovery};
    \node[stage, minimum width=25mm, right=of periph] (output) {\textbf{Output}\\classifier\\prediction};

    \draw[arrow] (input.east) -- (front.west);
    \draw[arrow, draw=teal!75!black] (front.east) -- ++(6mm,0) |- (analogpath.west);
    \draw[arrow] (front.east) -- ++(6mm,0) |- (digitalpath.west);
    \draw[arrow, draw=teal!75!black] (analogpath.east) -- (periph.west |- analogpath.east);
    \draw[arrow] (digitalpath.east) -- (periph.west |- digitalpath.east);
    \draw[arrow] (periph.east) -- (output.west);

    \end{tikzpicture}
    }
    \caption{\textbf{Hybrid organic-CIM inference stack.} Dense weight-stationary operators are executed on analog crossbars, whereas dynamic attention and control-heavy operators remain in the digital path.}
    \label{fig:system-architecture}
\end{figure}

\begin{table}[t]
\centering
\caption{\textbf{Analog/digital operator mapping rationale.} The split maximizes array utilization for weight-stationary operations while preserving precision for dynamic, irregular computations.}
\label{tab:operator-mapping}
\small
\begin{tabular}{@{}ll>{\raggedright\arraybackslash}p{6.5cm}@{}}
\toprule
\textbf{Operator} & \textbf{Domain} & \textbf{Rationale} \\
\midrule
Patch embedding (conv) & Analog & High weight reuse, regular access pattern; natural fit for crossbar MVM \\
QKV projections & Analog & Dense matrix-vector multiplications; weights stationary during inference \\
Attention output proj. & Analog & Post-attention linear projection; same characteristics as QKV \\
Feed-forward networks & Analog & MLP layers dominate parameter count; efficient crossbar utilization \\
\midrule
Depthwise convolutions & Digital & Low weight reuse, irregular spatial access; poor crossbar utilization \citep{wu2023bwq} \\
Layer normalization & Digital & Requires mean/variance computation across features; hard to parallelize on crossbar \\
$\mathbf{QK}^{\top}$ (attention scores) & Digital & Dynamic, input-dependent computation; irregular memory access \citep{lin2024hardsea} \\
Softmax & Digital & Non-linear transcendental function; precision-critical; no efficient analog implementation \citep{bettayeb2024memristorattention} \\
\bottomrule
\end{tabular}
\end{table}

\subsection{Modeling Physical Non-Idealities}
\label{subsec:modeling-nonidealities}

The framework injects quantization, cycle-to-cycle (C2C) and device-to-device (D2D) variability, ADC distortion, retention drift, and nonlinear-write surrogates directly into the forward path. Each analog-mapped weight $W_{ij}$ is converted into a differential conductance representation quantized to $n_{\text{states}}$ levels. We first decompose the weight into positive and negative branches:
\begin{equation}
    W_{ij}^{+} = \max(0, W_{ij}), \quad W_{ij}^{-} = \max(0, -W_{ij})
\end{equation}
These are then mapped to the physical conductance window $[G_{\min}, G_{\max}]$:
\begin{equation}
    G_{pos,ij} = G_{\min} + \frac{W_{ij}^{+}}{\max|W|} (G_{\max} - G_{\min})
\end{equation}
with $G_{neg,ij}$ calculated symmetrically. The differential-pair scheme enables signed-weight representation while partially suppressing common-mode noise \citep{wei2020voltagedifferential,kim2024sttmram}.
 Hardware-aware training (HAT) is implemented by keeping these forward perturbations active during optimization while propagating gradients through a straight-through estimator. For transformer deployments, the recovered post-array weight scale is modeled as an ideal calibrated digital rescaling step with fixed layer-wise factors, and its control overhead is therefore not itemized separately in the energy model.

First-order modeling of nonlinear write (\texttt{NL\_LTP/NL\_LTD}) and double-exponential retention is supported, with parameters anchored to measured DNTT transients \citep{vincze2026dualplasticity}. For optoelectronic sensing, we incorporate a physical frontend model (V6) characterized by the photoresponse exponent $\gamma_{phys}$, representing the power-law relation between light intensity and conductance, and the dark current $I_{dark}$, which sets the dynamic range floor. The released profiler uses representative edge-node energy and latency parameters.


\begin{figure}[ht]
    \centering
    \resizebox{0.98\textwidth}{!}{
    \begin{tikzpicture}[
        >=Latex,
        font=\small\rmfamily\upshape,
        box/.style={draw=black!60, rounded corners=4pt, thick, fill=black!3, minimum height=18mm, minimum width=28mm, align=center},
        accent/.style={draw=teal!70!black, fill=teal!10},
        arrow/.style={->, thick, draw=black!65}
    ]
    \node[box] (fp) {\textbf{FP32 weights}\\tensor $W$};
    \node[box, right=7mm of fp] (split) {\textbf{Branch split}\\$W_{ij}^{+}$ and $W_{ij}^{-}$};
    \node[box, accent, right=7mm of split] (map) {\textbf{Conductance map}\\$[G_{\min}, G_{\max}]$\\$n_{\mathrm{states}}$ levels};
    \node[box, accent, right=7mm of map] (profile) {\textbf{Profile effects}\\D2D / C2C\\retention / NL};
    \node[box, right=7mm of profile, minimum width=40mm] (readout) {\textbf{Readout + calibration}\\$G_{\mathrm{eff}} = G_{pos} - G_{neg}$\\$\hat{W}$};

    \draw[arrow] (fp.east) -- (split.west);
    \draw[arrow] (split.east) -- (map.west);
    \draw[arrow] (map.east) -- (profile.west);
    \draw[arrow] (profile.east) -- (readout.west);

    \end{tikzpicture}
    }
    \caption{\textbf{Behavioral weight-to-conductance mapping.} Floating-point weights are split into differential branches, mapped to a bounded conductance window, quantized, and perturbed by the selected device profile before differential readout.}
    \label{fig:weight-mapping}
\end{figure}

\subsection{Profile-Driven Substitution Interface}
\label{subsec:profile-interface}

The profile interface allows technology-specific parameters $(\sigma_{\text{D2D}}, \sigma_{\text{C2C}}, n_{\text{states}}, G_{\max}/G_{\min})$ to be substituted directly into the simulation workflow. In practice, this means that the same framework can absorb literature priors at present and measured-device characterizations later without changing the underlying evaluation pipeline. We therefore archive the optimization settings and random seeds needed for independent reruns. The resulting methodology is intended as a first-order behavioral framework for deployment-facing risk assessment, rather than as a pulse-accurate emulator.

\FloatBarrier
